\documentclass[../../../include/open-logic-section]{subfiles}

\begin{document}

\olfileid{sth}{replacement}{intro}
\olsection{جان پہچان}

بدل رکھن دا اصول اوہ مسلّمی خاکہ اے جیہڑا $\ZF$ تے~$\Z$ وچ فرق
پاندا اے۔ اسیں \crefrange{sth:ordinals::chap}{sth:spine::chap} وچ
ایہنوں لگاتار ورتیا اے۔ ایس باب وچ اسیں آخر ایہ سوال ویکھاں گے کہ
کی بدل رکھن دے اصول دا جواز اے؟

سوال نوں ہور صاف کرن لئی ایہ دھیان دین جوگی گل اے کہ بدل رکھن دا اصول
واقعی کافی \emph{طاقتور} اے۔ ایس باب وچ سانوں اندازہ ہووے گا کہ ایہ کِنّا
طاقتور اے (تے فیر \olref[sth][card-arithmetic][fix]{sec} وچ وی)۔
ایس توں سُجھدا اے کہ جواز دی واقعی لوڑ اے۔

اسیں دو قسماں دے جواز دی گل کراں گے۔ موٹی گل آکھیے تاں \emph{باہرلا}
جواز کسے مسلّمے نوں اوہدے حاصلاں راہیں جائز وکھاؤن دی کوشش اے۔
\emph{اندرلا} جواز ایہ وکھاؤن دی کوشش اے کہ زیر بحث ریاضیاتی تصور آپ
مسلّمے دی تائید کردے نیں۔ ایس باب وچ سانوں ایس گل دی ہور سمجھ آوے گی،
پر ایہ بس اک وڈی بحث دی شروعات اے۔ ہور لئی خاص کر کے
\citeauthor{Maddy1988a} (\citeyear{Maddy1988a} تے \citeyear{Maddy1988b})
ویکھو۔

\end{document}
